% SPDX-License-Identifier: Apache-2.0
% covers: Plic, Plic1, Plic2, Plic3, Plic4, Plic5, Plic6, Plic7, Plic8
\datasheet{Plic}{The platform-level interrupt controller: one to thirty-one sources and two targets, machine and supervisor, at the RISC-V PLIC's offsets.}

\dsfacts{\code{lib/parts/src/plic.rs}}%
{\code{txhdl\_parts::plic::Plic<N, EDGE>}, and \code{Plic1} to \code{Plic8}, which name the counts}%
{\code{//docs:vreteno}}

\subsection*{Function}

A controller gathers the interrupt lines of several sources into a
line for each of two targets, a hart's machine mode and its supervisor
mode (issue 1013), and lets software ask which source to serve. Each
source \textit{i}, numbered from 1, is an input line,
\code{srcs\_}\textit{i}\code{-1}, of the array of ports \code{srcs}. The
output \code{irq} is the machine target's line and \code{sirq} the
supervisor target's, \code{mip.SEIP}'s. The two targets share the
sources, their gateways and the pending bits; each has its own enable
bits, threshold and line. Software reaches the controller as an AXI-Lite
peripheral whose link is one port, \code{bus}, a \code{LitePort} of
the five channels: reads come on \code{bus\_ar} and are answered on
\code{bus\_r}, and writes come on \code{bus\_aw} and \code{bus\_w}
and are answered on \code{bus\_b}. It sits behind a
\code{LiteBridge<1, ..>}, which sends it only its own range.
The input \code{rst} clears the requests, both lines, the priorities,
and both targets' enable bits and thresholds.

It is one unit for any count of sources, \code{N} a const parameter
(issue 500): the sources are an array of ports and the priorities an
array of registers, \code{prio\_0} for source 1 onward, and the
lowering unrolls the loops over them for each \code{N}. The pending,
enable, active, held and previous-line words are 32 bits whatever
\code{N} is, bit \textit{i} for source \textit{i}, and the bits above
\code{N} stay zero; a source's number is 5 bits. \code{Plic1} to
\code{Plic8} are type aliases that name counts. \code{plic.rs} also names the
offsets, \code{PENDING}, \code{ENABLE}, \code{THRESHOLD} and
\code{CLAIM} for the machine target, \code{S\_ENABLE},
\code{S\_THRESHOLD} and \code{S\_CLAIM} for the supervisor target, and
\code{priority(i)}.

\subsection*{Parameters}

\begin{center}\footnotesize
\begin{tabular}{@{}lp{0.7\textwidth}@{}}
\toprule
Parameter & Meaning \\
\midrule
\code{N} & The count of sources, from 1 to 31, the most one word of
pending bits holds. Nothing checks it. \\
\code{EDGE} & A bit per source: bit \textit{i} set makes source
\textit{i} ask on a rising edge rather than while its line is high.
Bit 0 and bits above \textit{N} are ignored. The tables use 0, as the
board does. \\
\bottomrule
\end{tabular}
\end{center}

\subsection*{Register map}

Offsets from the controller's base, which are the RISC-V PLIC's for
two targets, contexts 0 and 1. The controller decodes the low 22 bits of an address. On
the board the base is \code{0x0c00\_0000}.

\begin{center}\footnotesize
\begin{tabular}{@{}llp{0.55\textwidth}@{}}
\toprule
Offset & Word & Access \\
\midrule
\code{0x4} $\times$ \textit{i} & Source \textit{i}'s priority, 3 bits &
Read and write, for \textit{i} from 1 to \textit{N}. \\
\code{0x1000} & The pending bits, bit \textit{i} for source \textit{i} &
Read only. \\
\code{0x2000} & The machine target's enable bits, the same way & Read
and write; bit 0 does not stick. \\
\code{0x2080} & The supervisor target's enable bits & The same. \\
\code{0x20\_0000} & The machine target's threshold, 3 bits & Read and
write. \\
\code{0x20\_0004} & The machine target's claim and complete & A read
claims; a write completes. \\
\code{0x20\_1000} & The supervisor target's threshold, 3 bits & Read
and write. \\
\code{0x20\_1004} & The supervisor target's claim and complete & The
same. \\
\bottomrule
\end{tabular}
\end{center}

Every other offset, offset 0 among them, reads as zero and ignores a
write.

\dstables{Plic}

\subsection*{Behaviour}

\begin{itemize}
\item \textbf{The gateway.} A level source asks while its line is high.
An edge source asks when its line is high and was low a cycle before
(\code{prev}), or when it holds an edge (\code{held}).
\item A request goes forward when its source is not \code{active}.
Going forward sets the source's bits in \code{pending} and in
\code{active}.
\item An edge that comes while its source is \code{active} sets its bit
in \code{held}, one deep, and it asks as soon as the source is no
longer active.
\item \textbf{The choice}, made for each target. A source is a
candidate for a target when it is pending, enabled in that target's
bits, \code{enable} or \code{senable}, and its priority is above that
target's threshold, \code{threshold} or \code{sthreshold}. The best
candidate has the highest priority, and the lowest number wins a tie.
A priority of 0 is never above a threshold, so it never interrupts.
\item \textbf{A claim.} A read of a target's claim word,
\code{0x20\_0004} or \code{0x20\_1004}, answers with that target's best
candidate's number, or 0 for none, and clears that source's
\code{pending} bit. The source stays \code{active}. So a source
enabled for both targets is served by whichever claims first, and the
other then finds it no longer pending. One read is taken a cycle, so
the two never claim together.
\item \textbf{A complete.} A write of a number from 0 to \textit{N} to
either claim word clears that source's \code{active} bit, whichever
target claimed it. A larger number changes nothing. A level source
whose line is still high then asks again at once.
\item \textbf{The lines.} \code{asserted} is set each cycle when the
machine target has a best candidate, and \code{sasserted} when the
supervisor target has one. \code{irq} follows \code{asserted} and
\code{sirq} follows \code{sasserted}, so each line is a register's
output and lags its choice by a cycle.
\item \textbf{The bus.} A read is answered in the cycle it is taken, and
a write is taken when its address and its word are both there, and
answered \code{Okay} at once. Every answer is \code{Okay}.
\end{itemize}

\subsection*{Verification}

\begin{itemize}
\item In \code{plic}, in \code{//lib/parts:parts\_test}, against a
\code{Plic3}:
\begin{itemize}
\item \code{the\_highest\_priority\_is\_claimed\_first\_and\_the\_lowest\_number\_wins\_a\_tie};
\item \code{a\_priority\_at\_or\_below\_the\_threshold\_does\_not\_interrupt};
\item \code{a\_claimed\_source\_asks\_again\_only\_after\_its\_complete};
\item \code{an\_edge\_source\_asks\_once\_per\_edge\_and\_a\_level\_source\_while\_high};
\item \code{the\_supervisor\_target\_is\_a\_target\_of\_its\_own}: its
enables and threshold read back, its line asks, and its claims and
completes follow the same rules;
\item \code{a\_source\_enabled\_for\_both\_targets\_is\_claimed\_once}:
both lines ask, the first claim takes it and the other target finds
none, and a complete by either ends the service.
\end{itemize}
\code{the\_eight\_source\_controller\_lowers} writes \code{Plic8}'s
Verilog.
\item \code{ex\_plic} runs \code{Plic3<0b100>} against a host that
claims and completes, last on the supervisor target, where source 3
enabled for it alone raises \code{sirq} and not \code{irq}. The build
simulates its netlist against that run under nvc and Verilator: \code{//docs:plic\_sim\_plic3\_tb\_test}
and \code{//docs:plic\_vsim\_test}.
\item \code{input\_comes\_by\_interrupt\_one\_byte\_each}, in
\code{//cpu/vreteno:board\_test}, runs the board's \code{Plic3} under a
compiled program that takes the serial port's input by interrupt.
\end{itemize}

\subsection*{Used by}

\code{Plic3<0>} in the board's design, \code{Board}, as its field
\code{plic} behind the bridge \code{pplic}: source 1 is the serial
port's receive interrupt, source 2 the board's \code{irq} input,
source 3 the Ethernet peripheral's, and
\code{irq} is the core's machine external interrupt, \code{mip.MEIP},
and \code{sirq} its supervisor external interrupt, \code{mip.SEIP}:
the core takes it as \code{seirq} and registers it beside MEIP
(issue 1094), which is how Linux in supervisor mode takes the serial
port's and the Ethernet port's interrupts. \code{Plic3<0b100>} in the example \code{ex\_plic}.

\subsection*{Limits}

It has two targets, contexts 0 and 1, and no others. Priorities are 3
bits, so from 0 to 7. A write ignores its strobe and writes the whole
field. The controller checks no address above bit 21, and relies on its
bridge for the rest. An edge source holds one edge during a service; a
second is lost.
